\documentclass[11pt,a4paper]{article}
\newcommand{\intestazione}{Dimero di spin --- 7. Dinamica sotto rumore}
% =====================================================================
%  Preambolo comune ai documenti sul dimero (serie "che cosa abbiamo fatto")
%  Incluso con % =====================================================================
%  Preambolo comune ai documenti sul dimero (serie "che cosa abbiamo fatto")
%  Incluso con % =====================================================================
%  Preambolo comune ai documenti sul dimero (serie "che cosa abbiamo fatto")
%  Incluso con \input{preambolo_dimero} subito dopo \documentclass.
% =====================================================================
\usepackage[T1]{fontenc}
\usepackage[utf8]{inputenc}
\usepackage[main=italian,provide=*]{babel}
\usepackage{amsmath,amssymb,mathtools}
\usepackage{graphicx}
\usepackage{booktabs}
\usepackage{colortbl}
\usepackage{xcolor}
\usepackage{geometry}
\usepackage{placeins}
\usepackage{caption}
\usepackage{enumitem}
\usepackage{fancyhdr}
\usepackage{needspace}
\usepackage{tikz}
\usetikzlibrary{arrows.meta,positioning,shapes.geometric,calc,fit,backgrounds}
\usepackage{hyperref}
\hypersetup{colorlinks=true,linkcolor=blu,citecolor=blu,urlcolor=blu}

\geometry{margin=2.5cm}

% --- colori coerenti con quelli delle figure (palette Okabe-Ito)
\definecolor{blu}{HTML}{0072B2}
\definecolor{arancio}{HTML}{E69F00}
\definecolor{verdes}{HTML}{009E73}
\definecolor{rossos}{HTML}{D55E00}
\definecolor{violas}{HTML}{CC79A7}
\definecolor{grigetto}{HTML}{E8E8E8}

\captionsetup{font=small,labelfont=bf,width=0.92\textwidth}

\pagestyle{fancy}
\fancyhf{}
\fancyhead[L]{\small\itshape\intestazione}
\fancyhead[R]{\small\thepage}
\renewcommand{\headrulewidth}{0.3pt}

% --- notazione
\newcommand{\ket}[1]{\lvert #1\rangle}
\newcommand{\bra}[1]{\langle #1 \rvert}
\newcommand{\media}[1]{\langle #1 \rangle}
\newcommand{\vs}[1]{\mathbf{s}_{#1}}

% --- riquadro "in una riga": il messaggio da portare a casa
\newcommand{\messaggio}[1]{%
  \begin{center}
  \begin{minipage}{0.93\textwidth}
  \setlength{\fboxsep}{9pt}%
  \colorbox{grigetto}{\begin{minipage}{\dimexpr\textwidth-2\fboxsep}
  \small\textbf{In una riga.}\enspace #1
  \end{minipage}}
  \end{minipage}
  \end{center}
}

% --- riquadro "come lo sappiamo": la verifica numerica dietro un'affermazione
\newcommand{\verifica}[1]{%
  \begin{center}
  \begin{minipage}{0.93\textwidth}
  \small\textcolor{verdes}{\rule{2.2em}{1.1pt}}\enspace
  \textcolor{verdes}{\textbf{Come lo sappiamo.}}\enspace #1
  \end{minipage}
  \end{center}
  \medskip
}
 subito dopo \documentclass.
% =====================================================================
\usepackage[T1]{fontenc}
\usepackage[utf8]{inputenc}
\usepackage[main=italian,provide=*]{babel}
\usepackage{amsmath,amssymb,mathtools}
\usepackage{graphicx}
\usepackage{booktabs}
\usepackage{colortbl}
\usepackage{xcolor}
\usepackage{geometry}
\usepackage{placeins}
\usepackage{caption}
\usepackage{enumitem}
\usepackage{fancyhdr}
\usepackage{needspace}
\usepackage{tikz}
\usetikzlibrary{arrows.meta,positioning,shapes.geometric,calc,fit,backgrounds}
\usepackage{hyperref}
\hypersetup{colorlinks=true,linkcolor=blu,citecolor=blu,urlcolor=blu}

\geometry{margin=2.5cm}

% --- colori coerenti con quelli delle figure (palette Okabe-Ito)
\definecolor{blu}{HTML}{0072B2}
\definecolor{arancio}{HTML}{E69F00}
\definecolor{verdes}{HTML}{009E73}
\definecolor{rossos}{HTML}{D55E00}
\definecolor{violas}{HTML}{CC79A7}
\definecolor{grigetto}{HTML}{E8E8E8}

\captionsetup{font=small,labelfont=bf,width=0.92\textwidth}

\pagestyle{fancy}
\fancyhf{}
\fancyhead[L]{\small\itshape\intestazione}
\fancyhead[R]{\small\thepage}
\renewcommand{\headrulewidth}{0.3pt}

% --- notazione
\newcommand{\ket}[1]{\lvert #1\rangle}
\newcommand{\bra}[1]{\langle #1 \rvert}
\newcommand{\media}[1]{\langle #1 \rangle}
\newcommand{\vs}[1]{\mathbf{s}_{#1}}

% --- riquadro "in una riga": il messaggio da portare a casa
\newcommand{\messaggio}[1]{%
  \begin{center}
  \begin{minipage}{0.93\textwidth}
  \setlength{\fboxsep}{9pt}%
  \colorbox{grigetto}{\begin{minipage}{\dimexpr\textwidth-2\fboxsep}
  \small\textbf{In una riga.}\enspace #1
  \end{minipage}}
  \end{minipage}
  \end{center}
}

% --- riquadro "come lo sappiamo": la verifica numerica dietro un'affermazione
\newcommand{\verifica}[1]{%
  \begin{center}
  \begin{minipage}{0.93\textwidth}
  \small\textcolor{verdes}{\rule{2.2em}{1.1pt}}\enspace
  \textcolor{verdes}{\textbf{Come lo sappiamo.}}\enspace #1
  \end{minipage}
  \end{center}
  \medskip
}
 subito dopo \documentclass.
% =====================================================================
\usepackage[T1]{fontenc}
\usepackage[utf8]{inputenc}
\usepackage[main=italian,provide=*]{babel}
\usepackage{amsmath,amssymb,mathtools}
\usepackage{graphicx}
\usepackage{booktabs}
\usepackage{colortbl}
\usepackage{xcolor}
\usepackage{geometry}
\usepackage{placeins}
\usepackage{caption}
\usepackage{enumitem}
\usepackage{fancyhdr}
\usepackage{needspace}
\usepackage{tikz}
\usetikzlibrary{arrows.meta,positioning,shapes.geometric,calc,fit,backgrounds}
\usepackage{hyperref}
\hypersetup{colorlinks=true,linkcolor=blu,citecolor=blu,urlcolor=blu}

\geometry{margin=2.5cm}

% --- colori coerenti con quelli delle figure (palette Okabe-Ito)
\definecolor{blu}{HTML}{0072B2}
\definecolor{arancio}{HTML}{E69F00}
\definecolor{verdes}{HTML}{009E73}
\definecolor{rossos}{HTML}{D55E00}
\definecolor{violas}{HTML}{CC79A7}
\definecolor{grigetto}{HTML}{E8E8E8}

\captionsetup{font=small,labelfont=bf,width=0.92\textwidth}

\pagestyle{fancy}
\fancyhf{}
\fancyhead[L]{\small\itshape\intestazione}
\fancyhead[R]{\small\thepage}
\renewcommand{\headrulewidth}{0.3pt}

% --- notazione
\newcommand{\ket}[1]{\lvert #1\rangle}
\newcommand{\bra}[1]{\langle #1 \rvert}
\newcommand{\media}[1]{\langle #1 \rangle}
\newcommand{\vs}[1]{\mathbf{s}_{#1}}

% --- riquadro "in una riga": il messaggio da portare a casa
\newcommand{\messaggio}[1]{%
  \begin{center}
  \begin{minipage}{0.93\textwidth}
  \setlength{\fboxsep}{9pt}%
  \colorbox{grigetto}{\begin{minipage}{\dimexpr\textwidth-2\fboxsep}
  \small\textbf{In una riga.}\enspace #1
  \end{minipage}}
  \end{minipage}
  \end{center}
}

% --- riquadro "come lo sappiamo": la verifica numerica dietro un'affermazione
\newcommand{\verifica}[1]{%
  \begin{center}
  \begin{minipage}{0.93\textwidth}
  \small\textcolor{verdes}{\rule{2.2em}{1.1pt}}\enspace
  \textcolor{verdes}{\textbf{Come lo sappiamo.}}\enspace #1
  \end{minipage}
  \end{center}
  \medskip
}


\title{\bfseries Più passi non è sempre meglio\\[2pt]
\large Documento 7 --- quantum simulation (Trotter) sotto rumore}
\author{Samuele Schiavi \\ \small Tesi triennale in Fisica, Università di Parma
--- Relatore: Prof.\ Alessandro Chiesa}
\date{}

\begin{document}
\maketitle
\thispagestyle{fancy}

\section{Il compromesso che il sistema chiuso non poteva mostrare}

Nel sistema chiuso (Documento 3), aumentare il numero di passi di Trotter $N$
migliorava sempre l'approssimazione dell'evoluzione $e^{-iHt}$: l'errore di
Trotter per passo è $O(\tau^2)$ con $\tau=t/N$, quindi l'errore totale scende
come $O(t^2/N)$, e a $N\to\infty$ si tende all'evoluzione esatta. Con il
rumore acceso questo non è più vero: ogni passo aggiunge gate, e ogni gate ha
una probabilità non nulla di sbagliare. Deve esistere un compromesso, e quindi
un $N$ ottimale $N^*$ --- questo documento lo trova.

L'Hamiltoniana si scompone come già nel Documento 3, $H=H_1+H_2$, con $H_1$ lo
scambio isotropo più il campo (fattorizza esattamente) e $H_2$ il solo
termine DM; un passo di Trotter al prim'ordine è
$U_\text{step}=e^{-iH_2\tau}e^{-iH_1\tau}$.

\section{Una trappola tecnica, già incontrata e già evitata}

Un breve richiamo di come è costruito il blocco, per non dover tornare al
Documento 3: il passo \emph{come si scriverebbe a mano}, gate per gate, userebbe
$5$ rotazioni a due qubit --- $R_{xx}(J\tau/2)$, $R_{yy}(J\tau/2)$,
$R_{zz}(J\tau/2)$ per $H_1$, più due $R_{zz}$ coniugate da Hadamard per $H_2$
--- esattamente come nel Documento 3. Qui però il passo non è costruito così:
la matrice $4\times4$ del passo, $U_\text{step}=e^{-iH_2\tau}e^{-iH_1\tau}$,
viene calcolata classicamente (\texttt{scipy.linalg.expm}) e passata a Qiskit
come un singolo \texttt{UnitaryGate} a 2 qubit, che il transpilatore
risintetizza direttamente nella base nativa $\{rz,sx,x,cx\}$ --- senza mai
passare dalle $5$ porte leggibili. Il risultato è lo stesso limite già
stabilito in Parte 1 per un operatore unitario a due qubit arbitrario (al più
$3$ CNOT), raggiunto qui per una via diversa ma equivalente: non il
compilatore che semplifica $5$ porte esplicite, ma la sintesi diretta di un
unitario dato in forma matriciale. Il singolo passo va transpilato
\textbf{una sola volta} al suo costo minimo (3 CNOT), e il blocco già
compilato va poi composto $N$ volte. Transpilare invece l'intero circuito
assemblato è un errore silenzioso: a un livello di ottimizzazione alto,
Qiskit riconosce che gli $N$ passi identici formano nel complesso un unico
unitario a 2 qubit, e lo risintetizza a costo \emph{costante}, indipendente
da $N$ --- cancellando esattamente la dipendenza da $N$ che è l'oggetto di
questo studio. La strategia di transpilazione a blocchi, adottata qui per la
prima volta sotto rumore, resta la regola per tutto il resto della pipeline.

\noindent\textit{Un dubbio legittimo, da chiarire subito}: sintetizzare il
passo da una matrice invece che da $5$ porte scritte a mano vuol dire che il
rumore di gate, qui, non vede un circuito vero? No. Dopo la transpilazione il
blocco \texttt{UnitaryGate} non esiste più: al suo posto c'è una sequenza
concreta di istruzioni --- per un passo, $8$ \texttt{rz}, $7$ \texttt{sx},
$3$ \texttt{cx}, ciascuna su un qubit preciso, verificato stampando
\texttt{qc.data} dopo la transpilazione (nessuna istruzione
\texttt{unitary} residua). È \emph{questa} sequenza, non la matrice di
partenza, a girare su \texttt{AerSimulator} con il \texttt{NoiseModel}
agganciato: il canale depolarizzante si applica per nome a ciascuna
istruzione, esattamente come farebbe su un circuito scritto a mano gate per
gate. Il rumore di gate resta quindi ``circuito vero'' in tutta questa
pipeline, qui come nel Documento 6. La sola scorciatoia dichiarata nella
Parte 2 riguarda il rumore di \emph{lettura}, non quello di gate --- vedi
Documento 8.

\section{L'osservabile: fedeltà rispetto allo stato bersaglio fisico}

\begin{equation}
\ket{\psi_\text{target}(t)} = e^{-iHt}\ket{\psi_0^\text{esatto}}, \qquad
F(N) = \bra{\psi_\text{target}(t)}\rho_N\ket{\psi_\text{target}(t)}
\end{equation}
dove $\rho_N$ è lo stato ottenuto da: preparazione VQE rumorosa (Documento 6)
$\to$ $N$ passi di Trotter rumorosi. Una singola quantità impacchetta tre
contributi diversi: la preparazione, costante in $N$; l'errore di Trotter, che
decresce con $N$; il rumore di gate accumulato, che cresce con $N$.

\noindent\textit{Un dubbio legittimo}: se $H_1$ e $H_2$ non commutano per
$D\neq0$, come si fa a calcolare $\ket{\psi_\text{target}(t)}$ \emph{senza}
Trotter? Non serve affatto la scomposizione $H=H_1+H_2$: $\ket{\psi_\text{
target}(t)}$ si ottiene da $U=\texttt{scipy.linalg.expm}(-iHt)$, applicato
alla matrice $4\times4$ dell'Hamiltoniana \textbf{intera} (DM incluso), non
ai suoi due pezzi separati. È un calcolo classico puro --- un'esponenziale
di matrice su un computer normale, non un circuito --- e funziona per
qualunque $H$ hermitiana, che i suoi eventuali pezzi commutino o no: la
domanda ``$H_1$ e $H_2$ commutano?'' ha senso solo quando si vuole
\emph{fattorizzare} l'esponenziale in pezzi implementabili come gate,
cosa che serve per costruire il circuito di $\rho_N$, non per calcolare
$\ket{\psi_\text{target}(t)}$. I due lati del confronto sono quindi ottenuti
con metodi deliberatamente indipendenti --- altrimenti il confronto
sarebbe circolare, Trotter contro Trotter.

\section{Prima di fidarsi della curva: un secondo cammino di calcolo}

A rumore nullo, l'evoluzione di Trotter si può calcolare anche per pura
potenza di matrice in \texttt{numpy}, senza passare da nessun oggetto
Qiskit --- un controllo indipendente contro un bug nella costruzione del
circuito (per esempio nell'ordine dei blocchi, o nel segno di $H_2$) che
produca comunque una curva plausibile ma sbagliata.

\begin{center}
\renewcommand{\arraystretch}{1.25}
\begin{tabular}{@{}lcc@{}}
\toprule
$N$ & Circuito Qiskit & Potenza di matrice (\texttt{numpy})\\
\midrule
$8$ & $0.9454850$ & $0.9454850$\\
$20$ & $0.9921835$ & $0.9921835$\\
\bottomrule
\end{tabular}
\end{center}

\verifica{Coincidenza esatta, alla precisione mostrata, fra due cammini di
calcolo completamente indipendenti --- uno passa da un circuito quantistico
simulato, l'altro da una moltiplicazione di matrici $4\times4$ pura. Rieseguito
anche in questa sessione: stessi valori, alla cifra.}

\section{Perché serve il rumore: confronto diretto}

\begin{center}
\renewcommand{\arraystretch}{1.15}
\small
\begin{tabular}{@{}cccc@{}}
\toprule
$N$ & $F$ (rumore nullo) & $F$ (\texttt{ibm\_torino}) & \\
\midrule
$1$ & $0.7799372$ & $0.7567291$ & \\
$2$ & $0.4776280$ & $0.4624084$ & \\
$5$ & $0.8323751$ & $0.7560090$ & \\
$8$ & $0.9454850$ & $0.8177249$ & $\approx N^*$\\
$10$ & $0.9666876$ & $0.8087205$ & \\
$20$ & $0.9921835$ & $0.7097565$ & \\
$40$ & $0.9980771$ & $0.5427382$ & \\
$80$ & $0.9995212$ & $0.3567462$ & \\
$160$ & $0.9998804$ & $0.2686760$ & \\
\bottomrule
\end{tabular}
\end{center}

Senza rumore, $F(N)\to1$ al crescere di $N$: puro errore di Trotter, in
diminuzione monotona per $N\gtrsim5$, nessun compromesso, nessun massimo. Con
il rumore di riferimento, la stessa curva smette di crescere oltre $N=8$ e
comincia a scendere: è \emph{solo} il rumore di gate accumulato a creare un
massimo finito --- senza rumore l'ottimo sarebbe $N\to\infty$.

\section{Risultato: la curva $F(N)$ e $N^*$}

\begin{figure}[htbp]
\centering
\includegraphics[width=0.72\textwidth]{figure/fig_passo3_F_vs_N}
\caption{Fedeltà $F(N)$ rispetto allo stato bersaglio fisico, in funzione del
numero di passi di Trotter, a rumore di riferimento (\texttt{ibm\_torino}).
Crescita rapida per $N$ piccolo (l'errore di Trotter domina), massimo a
$N^*=8$, decrescita lenta per $N$ grande (il rumore di gate accumulato
domina).}
\end{figure}

\FloatBarrier

\begin{center}
\renewcommand{\arraystretch}{1.25}
\begin{tabular}{@{}lc@{}}
\toprule
Quantità & Valore\\
\midrule
$N^*$ & $8$\\
$F(N^*)$ & $0.817725$\\
\bottomrule
\end{tabular}
\end{center}

\section{Un dettaglio che una prima occhiata avrebbe potuto perdere}

\begin{figure}[htbp]
\centering
\includegraphics[width=0.68\textwidth]{figure/fig_passo3_zoom_minimo}
\caption{Zoom sulla regione non perturbativa, $N=1$ a $8$. Il minimo (non il
semplice calo che ci si aspetterebbe) è a $N=3$, identico a rumore nullo e a
rumore acceso --- prova diretta che è un effetto di Trotter puro, non un
artefatto del rumore o della simulazione.}
\end{figure}

\FloatBarrier

La curva non è monotona per $N$ piccolo: parte relativamente alta a $N=1$
($F\approx0.78$), scende in modo sostanziale --- la fedeltà quasi si dimezza
--- fino a un minimo a $N=3$ ($F\approx0.40$), poi risale con continuità verso
il massimo. Non è rumore numerico: lo stesso minimo, nella stessa posizione,
compare identico anche a rumore nullo --- è il regime non ancora perturbativo
dell'errore di Trotter ($\tau=t/N=2,\,1,\,0.67$ per $N=1,2,3$: tutt'altro che
piccolo), non un artefatto della simulazione.

\verifica{Rieseguito il calcolo per $N=3$ e $N=8$ separatamente da tutto il
resto della tabella: $N=3$ dà $F=0.4018$ (rumore nullo) e $F=0.3892$
(rumoroso), $N=8$ dà $F=0.9455$ e $F=0.8177$ --- entrambi coerenti a tutte le
cifre con la tabella sopra. Il minimo a $N=3$ non è un artefatto del modo in
cui è stata costruita la griglia di punti.}

\section{Cosa non entra in questa fedeltà}

Il readout non compare in questa metrica: è una fedeltà fra stati (operatore
densità contro stato puro), non una misura di un osservabile via ancilla. Il
readout entrerà per la prima volta al documento successivo, dove la quantità
osservata passa attraverso una misura reale.

\messaggio{Il compromesso Trotter/rumore trovato qui ($N^*=8$) non è
universale: dipende dal tempo $t$, dal punto di lavoro fisico, e dal livello
di rumore --- lo si vedrà esplicitamente nello scan sui parametri di rumore
(Documento 9). Il documento successivo (correlatori dinamici) userà lo stesso
principio --- Trotter per l'evoluzione, transpilazione a blocchi mai
ricomposta per intero --- ma su un circuito diverso, un test di Hadamard con
ancilla, e troverà un $N^*$ diverso ($N^*=5$, non $8$): la posizione del
compromesso dipende dal circuito specifico, un fatto da verificare di nuovo
lì, non da assumere per analogia da qui.}

\end{document}
